\section{Where real campaigns sit, and what passive monitoring cannot do}
\label{sec:world}

\paragraph{Documented campaigns.} The campaigns reported in 2026 occupy two corners, neither of them the covert, self-written corner we studied most. At one are high-volume campaigns: the largest Anthropic has reported ran more than 151 million exchanges in three months through 3{,}500--5{,}000 accounts, several hundred exchanges per account per day, and its accounts were linked because they all used one fixed prompt \citep{anthropic2026threat}. Per-account volume monitors and template coordination, which our detectors caught almost every week, see such campaigns. At the other corner are \emph{relays}: one forwarded almost 300{,}000 customer requests in ten days through 5{,}380 accounts, about six per account per day \citep{anthropic2026threat}.

\paragraph{Why relays are the hard case.} Our covert campaigns were limited by supply: an attacker without users could not obtain or write enough prompts that are both useful and user-like. A relay has an unlimited supply. If a relaying account's traffic has the same joint law as the users it relays, Proposition~\ref{prop:law} leaves no query-only monitor able to separate it. That ``if'' is an assumption: a relay that interleaves many users' sessions in one account, or forwards only high-value requests, may differ in session coherence or topicality. We did not simulate relays; it is the most important open experiment.

\paragraph{The provider's own randomness.} Every threshold here is learned from history, which drifts and can be poisoned. A provider has one resource no attacker controls: if it throttles a flagged group in randomly chosen time blocks, it can test whether the work fails over to other accounts, with a null that follows from the randomisation alone, whatever benign traffic does (Fisher's randomisation test in a switchback design; \citealp{fisher1935design,bojinov2023design}; Proposition~\ref{prop:rand} in Appendix~\ref{sec:world-full}). In simulation on WildChat populations the null held on real, drifting traffic, and campaigns that reassigned throttled work were caught in every run. But a benign team with legitimate failover fired the test 82--98\% of the time, throttling harms legitimate accounts in a flagged group, and twelve weekly blocks take a quarter of a year (Appendix~\ref{app:randomized}).
